\scanpage{19}
\uline{4} Mesures spectrales et réarrangements spectraux d'opérateurs.

Dans toute la suite, $\mathcal A$ est une algèbre de von Neumann (algèbre auto-adjointe \unsure{fblt} fermée \deleted{\unsure{disquée}} avec unité d'opérateurs dans un espace de Hilbert). \inserted{Une trace} \deleted{\unsure{une trace} trace normale (} sur $\mathcal A$ est une \deleted{\unsure{fonction}} fonction $T$ sur la partie positive $\mathcal A^+$ de $\mathcal A$, à valeurs $0\leq T(A)\leq+\infty$, \deleted{additive, \unsure{positivement homogène}} telle que $T(A+B)=T(A)+T(B)$ et $T(\lambda A)=\lambda T(A)$ ($A,B\in\mathcal A^+$, $\lambda$ \deleted{$\unsure{>0}$} scalaire $\geq0$; on pose $0.(+\infty)=0$), et invariante par transformations unitaires. Alors l'ens. des $A\in\mathcal A^+$ tels que $T(A)<\infty$ est la partie positive $\underline a^+$ d'un idéal bilatère \inserted{auto-adjoint} $\underline a$, \inserted{l'idéal de définition de $T$,} et $T$ se prolonge en une forme linéaire sur $\underline a$, notée encore $T$. On a
\[
T(AB)=T(BA)\qquad A\in\underline a,\ B\in\mathcal A
\]
$T$ est dite \deleted{\unsure{essentielle}} semi-finie si pour tout \deleted{$A\in\mathcal A^+$, $A\neq0$, \unsure{existe} $B$} $A\in\mathcal A$, $A>0$, existe $B$ tel que $0<B<A$ et $T(B)\unsure{<|\leq}+\infty$. $T$ est dite normale si
\[
T(\mathop{\rm Sup} A_i)=\mathop{\rm Sup} T(A_i)
\]
pour toute famille filtrante croissante majorée \inserted{$(A_i)$} dans $\mathcal A^+$. Si $T$ est normale, $T$ est \deleted{finie} \unsure{s. finie|semi-finie} si et seulement si \deleted{\unsure{de}} son idéal de définition $\underline a$ est \uline{\unsure{fblt} dense} dans $\mathcal A$.

\inserted{$T$ est dite fidèle si \deleted{$\unsure{A\in\mathcal A^+}$} $A\geq0$, $T(A)=0\Longrightarrow A=0$.}

Pour la suite, \uline{on suppose $\mathcal A$ muni une fois pour toutes d'une trace normale \unsure{s. finie|semi-finie} \deleted{$\unsure{T}$} \inserted{fidèle $T$}} dont l'idéal de définition est $\underline a$. Soit $\underline b$ l'adhérence de $\underline a$ pour la topologie de la norme, c'est un idéal bilatère auto-adjoint \unsure{uniformt} fermé.
\mspage{20}
